\documentclass[fleqn]{article}
\usepackage{amsmath, amssymb, amsthm}
\usepackage{fullpage}
\usepackage{listings}
\usepackage{stix}
\usepackage[utf8]{inputenc}
\lstset{ basicstyle=\ttfamily, }

\begin{document}

\section*{Excercise 1.5 1}

$Q(x, y, z)$

$y$ is in the domain of the function computed by the program $x$ if and only if there is a $z$ that makes $Q$ true.

Construct $Q$ in a way that it is try when $z$ causes (or encodes) some notion of $y \in E_x$

$\phi_x(?z?) = {y}$

$\phi_x^{-1}(y) = (z)_1$

$Q$ says whether $(z)_1$ is the inverse of y.

$y \in E_x$ \& $Q(x, y, z)$ $\implies$ $\phi_x((z)_1) = y$

$Q$ checks that $\phi_x((z)_1) = y$ and $\phi_x^{-1}(y) = (z)_1$

\bigskip

This is the first step of constructing an inverse function that "sweeps" the range space to find the domain.

\section*{3.2 1}

Construct computable function $k(e)$ that encodes not $M(e)$.

$M_x\simeq\phi_e(x)$

$\neg M_x \simeq \phi_k(e)(x)$

$k(e) \equiv e, 1 \dot{-} R_1$

\bigskip

Using the concatenation construction of programs and the negation program.

\section*{3.2 2}

$k(x): E_{k(x)} = W_x$

Total computable function that the domain of the function on a program is the range of the program.

$k(x) \equiv \phi_x^{-1}$

\bigskip

$k(x)$ computes the inverse function for program with number $x$, see 1.5 for such a construction.

\section*{3.2 3}

$S(x, y): E_{s(x, y)} = E_x\cup E_y$

A total computable function $S(x, y)$ whose domain is the union of domains of both programs $x$ and $y$.

$$
\phi_{S(x, y)}(z) =
\begin{cases}
    \phi_x(\frac{z}{2}) \text{ if $z$ even} \\
    \phi_y(\frac{z - 1}{2}) \text{ if $z$ odd}
\end{cases}
$$

\bigskip

$S(x, y)$ uses even/odd encoding to partition the naturals into the program $x$ case and program $y$ case.

\end{document}